\begin{resumo}
    Este trabalho apresenta as etapas do processo de desenvolvimento do
    Solidarize, uma plataforma de ações sociais e voluntariado que tem como objetivo
    facilitar a conexão entre Organizações da Sociedade Civil (OSCs) e pessoas
    interessadas em colaborar por meio do voluntariado. A plataforma permite que as
    OSCs publiquem oportunidades de voluntariado e que os voluntários encontrem
    vagas alinhadas aos seus interesses e disponibilidade. Além disso, a aplicação
    possibilita ao voluntário o compartilhamento de experiências em projetos que já
    participou. Serão apresentados no documento as seguintes etapas: área de
    abrangência, apresentando uma breve contextualização sobre o terceiro setor e
    suas organizações na sociedade civil; metodologia da pesquisa com o resultado da
    pesquisa de levantamento dos requisitos junto ao público-alvo da plataforma, além
    de uma análise comparativa de soluções similares ao proposto no desenvolvimento
    deste trabalho; os requisitos inicialmente definidos para o desenvolvimento do
    projeto e por fim o processo de modelagem apresentando os modelos que irão
    representar diferentes perspectivas da plataforma.
    \vspace{0.5cm}

    \textbf{Palavras-chave:} Voluntariado; Organizações da Sociedade Civil; Causas
    Sociais.
\end{resumo}